(1) 因 \((\boldsymbol\sigma\cdot\boldsymbol n)^2=I\) 且 \(|\boldsymbol n|=1\)，本征值为 \(\pm1\)，对应自旋沿或反向于 \(\boldsymbol n\) 的态。哈密顿量本征能量为 \(E_+=-\hbar\Omega/2\)（沿场低能）及 \(E_-=\hbar\Omega/2\)（反向高能），始终具有固定能隙 \(\hbar\Omega\)。这为绝热跟随提供避免能级交叉的条件；快速场变化时仍可能发生非绝热跃迁。

(2) 选择北极规范的沿场本征旋量 \(|+\boldsymbol n\rangle=(\cos(\theta/2),e^{i\phi}\sin(\theta/2))^T\)，\(\phi=\omega t\)。其 Berry 连接 \(i\langle+|\partial_\phi+\rangle=-\sin^2(\theta/2)\)，故一周几何相位 \(\gamma_+=-2\pi\sin^2(\theta/2)=-\pi(1-\cos\theta)\)（模 \(2\pi\)）。另一带的几何相位符号相反 \(\gamma_-=\pi(1-\cos\theta)\)，对应 Bloch 球闭合轨迹所围立体角 \(2\pi(1-\cos\theta)\)。

(3) 动力学相位为 \(\delta_\pm=-(1/\hbar)\int_0^{\mathcal T}E_\pm dt\)，因此 \(\delta_+=+\Omega\mathcal T/2,\ \delta_-=-\Omega\mathcal T/2\)。总相位分别 \(\alpha_+=\Omega\mathcal T/2-\pi(1-\cos\theta)\)，\(\alpha_-=-\Omega\mathcal T/2+\pi(1-\cos\theta)\)，理解为相位模 \(2\pi\)。几何项不依赖具体缓慢角速度，而动力学项与持续时间成正比。

(4) 若初态为 \((|+\boldsymbol n(0)\rangle+|-\boldsymbol n(0)\rangle)/\sqrt2\)，一周后两分量相对相位变化 \(\Delta\alpha=\alpha_+-\alpha_-=\Omega\mathcal T-2\pi(1-\cos\theta)\)。干涉测量可通过控制旋转方向将几何项与动力学项分离。\(\theta\to0\) 时立体角趋零，Berry 相位消失。绝热条件除了 \(\omega\sin\theta\ll\Omega\) 还要求相关控制噪声慢且能隙不闭合；否则不能用单纯相位跟随描述真实演化。
